\scanpage{29}
% archive sheet number: 28
\mspage{-39-}\hfill n$^{\mathrm{o}}$ 272

Corollaire 1.- \uline{Pour toute famille} $a_i\handwritten{\in} A_i$ $(1\handwritten{\leqslant}i\handwritten{\leqslant}n)$, \uline{il existe} $x\handwritten{\in}B$ \uline{tel que} $x\handwritten{\equiv}a_i\,(\mathrm{mod.}\ \underline{p}_i)$.

Puisque les $\underline{p}_i$ sont des idéaux maximaux de $B$, on a $B_{\underline{p}_i}\handwritten{/}\underline{p}_i B_{\underline{p}_i}=B\handwritten{/}\underline{p}_i$, et l'on peut donc supposer que $a_i\handwritten{\in}B$ pour tout $i$. Dans ce cas, le corollaire signifie que l'application $B\handwritten{\to}\handwritten{\prod} B\handwritten{/}\underline{p}_i$ est surjective, ce qui est un résultat bien connu (cf.Ch.I, \handwritten{\S}2, cor. à la prop.4).

Corollaire 2.- \uline{Il existe des éléments} $x_i\handwritten{\in}K$ \uline{tels que} $v_i(x_i)=0$ \uline{et} $v_j(x_i)\handwritten{>}0$ \uline{pour} $i\handwritten{\ne}j$.

Pour chaque indice $i$, on applique le cor.1 à la famille $a_j=0$ si $j\handwritten{\ne}i$, $a_i=1$.

Corollaire 3.- \uline{Tout anneau de valuation contenant} $B$ \uline{contient l'un des} $A_i$.

Soit $A$ un anneau de valuation contenant $B$, et soit $\underline{p}=\underline{m}(A)\handwritten{\cap}B$; il existe un idéal maximal $\underline{p}_i$ de $B$ contenant $\underline{p}$, d'ou $A_i=B_{\underline{p}_i}\handwritten{\unsure{<|\subset}}B_{\underline p}\handwritten{\subset}A$. \handwritten{[}En fait, on a $A=B_{\underline p}$ en vertu de la prop.1\unsure{;|,} mais on s'en fout.\handwritten{]}

2. \uline{Valuations indépendantes}.

Définition 1.- \uline{Soient} $A$ \uline{et} $A'$ \uline{deux anneaux de valuation d'un même corps} $K$. \uline{On dit que} $A$ \uline{et} $A'$ \uline{sont indépendants s'il n'existe aucun sous-anneau de} $K$, \uline{distinct de} $K$, \uline{les contenant}.

Bien entendu, on dit que deux valuations $v$ et $v'$ de $K$ sont \uline{indépendantes} si leurs anneaux sont indépendants.

\uline{Exemple}. Si $A$ et $A'$ sont \uline{distincts} et sont tous deux des anneaux de valuation de \uline{rang} 1, ils sont indépendants, en vertu de la prop.7 du \handwritten{\S}4.

Proposition 3.- (\uline{théorème d'approximation}). \uline{Soient} $v_i$ $(1\handwritten{\leqslant}i\handwritten{\leqslant}n)$ \uline{des valuations deux à deux indépendantes d'un corps} $K$, \uline{et soient} $\handwritten{\Gamma}_i$ \uline{leurs groupes des ordres}. \uline{Soient} $a_i\handwritten{\in}K$ \uline{et} $\handwritten{\gamma}_i\handwritten{\in}\handwritten{\Gamma}_i$ $(1\handwritten{\leqslant}i\handwritten{\leqslant}n)$. \uline{Il existe alors} $x\handwritten{\in}K$ \uline{tel que} $v_i(x-a_i)\handwritten{\geqslant}\handwritten{\gamma}_i$ \uline{pour tout} $i$.

